
\documentclass{article}
\usepackage{colortbl}
\usepackage{makecell}
\usepackage{multirow}
\usepackage{supertabular}

\begin{document}

\newcounter{utterance}

\centering \large Interaction Transcript for game `eqbench', experiment `eq\_bench', episode 71 with reflexion{-}llama{-}3.1{-}8b{-}instruct.
\vspace{24pt}

{ \footnotesize  \setcounter{utterance}{1}
\setlength{\tabcolsep}{0pt}
\begin{supertabular}{c@{$\;$}|p{.15\linewidth}@{}p{.15\linewidth}p{.15\linewidth}p{.15\linewidth}p{.15\linewidth}p{.15\linewidth}}
    \# & \multicolumn{2}{c}{Player} && \multicolumn{2}{c}{Game Master} \\
    \hline

    \theutterance \stepcounter{utterance}  
    & & & \multicolumn{4}{p{0.6\linewidth}}{
        \cellcolor[rgb]{0.9,0.9,0.9}{
            \makecell[{{p{\linewidth}}}]{
                \texttt{\tiny{[P1$\langle$GM]}}
                \texttt{You have attempted this same task before and did not succeed. These are the plans you wrote after those attempts.} \\
\texttt{Trial \#1: New plan: I got the emotions right, but my scores were not accurate. Upon re{-}reading the dialogue, I realize that Richard is trying to manipulate Emily into publishing the information, and she is not buying it. He's being evasive and using euphemisms, which suggests that he's not confident in his story. He's also trying to pressure her into publishing without verifying the evidence, which makes him seem impatient. Emily, on the other hand, is skeptical and not easily swayed, which doesn't suggest that she's anxious or ashamed. } \\
\\ 
\texttt{To improve my scores, I will pay closer attention to the tone and language used in the dialogue, and try to identify the underlying emotions and motivations of the characters. I will also consider the context and the power dynamics at play in the conversation. } \\
\\ 
\texttt{Impatient: 9} \\
\texttt{Confident: 2} \\
\texttt{Anxious: 0} \\
\texttt{Ashamed: 0} \\
\\ 
\texttt{Trial \#2: New plan: Upon re{-}reading the dialogue, I realize that I was correct in identifying Richard's impatience and lack of confidence, but I didn't consider the emotional impact of Emily's skepticism and Richard's evasiveness on his emotions. Emily's sharp responses and Richard's defensive tone suggest that he's not only impatient but also anxious about being caught in his manipulation. His use of euphemisms and evasive language also implies that he's trying to hide something, which could make him feel ashamed. } \\
\\ 
\texttt{To improve my scores, I will pay closer attention to the emotional cues and power dynamics in the conversation, and consider how the characters' emotions and motivations are intertwined. I will also look for subtle hints and contradictions in the dialogue that could reveal more about the characters' emotional states.} \\
\\ 
\texttt{Impatient: 9} \\
\texttt{Confident: 2} \\
\texttt{Anxious: 6} \\
\texttt{Ashamed: 4} \\
\\ 
\\ 
\\ 
\texttt{Your task is to predict the likely emotional responses of a character in this dialogue:} \\
\\ 
\texttt{Richard: Emily, always a pleasure. You're looking... tenacious as always.} \\
\texttt{Emily: Cut the flattery, Richard. What do you want?} \\
\texttt{Richard: Straight to the point. I like that. I have a story for you. A big one.} \\
\texttt{Emily: I'm listening.} \\
\texttt{Richard: It's about my opponent. There are some... unsavory details about his past that the public deserves to know.} \\
\texttt{Emily: And you're sharing this out of the goodness of your heart?} \\
\texttt{Richard: Let's just say, I believe in transparency.} \\
\texttt{Emily: You mean manipulation.} \\
\texttt{Richard: That's a harsh word, Emily. I prefer 'strategic communication'.} \\
\texttt{Emily: Call it what you want, Richard. I'm not your puppet.} \\
\texttt{Richard: Of course not. I respect your work, Emily. That's why I came to you.} \\
\texttt{Emily: I'm sure.} \\
\texttt{Richard: I have evidence. It's in this envelope.} \\
\texttt{Emily: And you expect me to just publish this without verifying?} \\
\texttt{Richard: I trust your judgment, Emily. You'll see it's all true.} \\
\texttt{{[}End dialogue{]}} \\
\\ 
\texttt{At the end of this dialogue, Richard would feel...} \\
\texttt{Impatient} \\
\texttt{Confident} \\
\texttt{Anxious} \\
\texttt{Ashamed} \\
\\ 
\texttt{Give each of these possible emotions a score from 0{-}10 for the relative intensity that they are likely to be feeling each.} \\
\\ 
\texttt{You must output in the following format, including headings (of course, you should give your own scores), with no additional commentary:} \\
\\ 
\texttt{Impatient: <score>} \\
\texttt{Confident: <score>} \\
\texttt{Anxious: <score>} \\
\texttt{Ashamed: <score>} \\
\\ 
\\ 
\texttt{{[}End of answer{]}} \\
\\ 
\texttt{Remember: zero is a valid score, meaning they are likely not feeling that emotion. You must score at least one emotion > 0.} \\
\\ 
\texttt{Your answer:} \\
            }
        }
    }
     \\ \\

    \theutterance \stepcounter{utterance}  
    & \multicolumn{4}{p{0.6\linewidth}}{
        \cellcolor[rgb]{0.561,0.737,0.561}{
            \makecell[{{p{\linewidth}}}]{
                \texttt{\tiny{[P1$\rangle$GM]}}
                \texttt{Impatient: 8} \\
\texttt{Confident: 3} \\
\texttt{Anxious: 7} \\
\texttt{Ashamed: 6} \\
            }
        }
    }
    & & \\ \\

    \theutterance \stepcounter{utterance}  
    & & & \multicolumn{2}{p{0.3\linewidth}}{
        \cellcolor[rgb]{0.9,0.9,0.9}{
            \makecell[{{p{\linewidth}}}]{
                \texttt{\tiny{[GM$|$GM]}}
                \texttt{Impatient: 8} \\
\texttt{Confident: 3} \\
\texttt{Anxious: 7} \\
\texttt{Ashamed: 6} \\
            }
        }
    }
    & & \\ \\

    \theutterance \stepcounter{utterance}  
    & & & \multicolumn{2}{p{0.3\linewidth}}{
        \cellcolor[rgb]{0.9,0.9,0.9}{
            \makecell[{{p{\linewidth}}}]{
                \texttt{\tiny{[GM$|$GM]}}
                \texttt{Impatient: 3} \\
\texttt{Confident: 8} \\
\texttt{Anxious: 2} \\
\texttt{Ashamed: 0} \\
            }
        }
    }
    & & \\ \\

    \theutterance \stepcounter{utterance}  
    & & & \multicolumn{2}{p{0.3\linewidth}}{
        \cellcolor[rgb]{0.9,0.9,0.9}{
            \makecell[{{p{\linewidth}}}]{
                \texttt{\tiny{[GM$|$GM]}}
                \texttt{game\_result = LOSE} \\
            }
        }
    }
    & & \\ \\

    \theutterance \stepcounter{utterance}  
    & & \multicolumn{4}{p{0.6\linewidth}}{
        \cellcolor[rgb]{0.95,0.95,0.95}{
            \makecell[{{p{\linewidth}}}]{
                \texttt{\tiny{[P1$\leftrightarrow$P]}}
                \texttt{Post{-}episode Reflexion (not sent to GM):} \\
\texttt{New plan: Upon re{-}reading the dialogue, I realize that I was correct in identifying Richard's impatience and lack of confidence, but I didn't consider the emotional impact of Emily's skepticism and Richard's evasiveness on his emotions. However, I also realize that I overestimated Richard's anxiety and shame. Emily's sharp responses and Richard's defensive tone suggest that he's not only impatient but also defensive, which could be a sign of anxiety. However, his use of euphemisms and evasive language also implies that he's trying to hide something, but it's not necessarily shame. } \\
\\ 
\texttt{To improve my scores, I will pay closer attention to the emotional cues and power dynamics in the conversation, and consider how the characters' emotions and motivations are intertwined. I will also look for subtle hints and contradictions in the dialogue that could reveal more about the characters' emotional states. I will also consider the possibility that Richard's anxiety is more related to his manipulation being exposed rather than shame. } \\
\\ 
\texttt{Impatient: 9} \\
\texttt{Confident: 2} \\
\texttt{Anxious: 8} \\
\texttt{Ashamed: 2} \\
            }
        }
    }
    & \\ \\

\end{supertabular}
}

\end{document}
